Sei $\ell$ die Länge und $m$ die Masse des Auslegers, $M$ die Masse der Last und $\alpha$ der Winkel zwischen Seil und Ausleger.

\begin{abcliste}
\abc
Nimm das Gelenk als Drehachse: Seine Kraft hat dann keinen Hebelarm. Die Gewichtskräfte drehen den Ausleger nach unten (die des Auslegers greift in seiner Mitte an), das Seil mit seiner Komponente $F_\mathrm{S}\sin\alpha$ senkrecht zum Ausleger nach oben:
\begin{align}
 F_\mathrm{S}\sin\alpha\cdot\ell &= m\,g\cdot\frac{\ell}{2} + M\,g\cdot\ell
\end{align}
$\ell$ kürzt sich weg: Die Länge des Auslegers spielt keine Rolle.
\begin{align}
 F_\mathrm{S} &= \FSF \\
   &= \frac{\left(\frac{\mAu}{2}+\ML\right)\cdot\g}{\sin\al} \\
   &= \FS \approx \result{\FSP{3}{2}}
\end{align}

\abc
Das Seil zieht den Ausleger mit der Komponente $F_\mathrm{S}\cos\alpha$ zur Wand hin; sonst wirkt keine horizontale Kraft ausser der des Gelenks, also drückt das Gelenk gleich stark zurück:
\begin{align}
 F_\mathrm{G} &= \FGF \\
   &= \FS\cdot\cos\al \\
   &= \FG \approx \result{\FGP{0}{2}}
\end{align}
\end{abcliste}
